%%%%%
%
\section{D\'ecalage}

\itemstyle{Op\'erateur~:} {\sf Decale}
\index{Decale}

%%%%%OBJET%%%%
\itemstyle{Objet~:}
D\'ecalage temporel d'un param\`etre par
interpolation lin\'eaire. 

%%%%%ARGUMENTS%%%%
\bigskip
%%%%%ENTREES%%%%
\begin{tabarg}
\multicolumn{3}{l}{\titretab {Entr\'ees}}\\
\hline
$X$ & Vecteur, {\tt float} & Param\`etre d'entr\'ee \\
$fe$ & Coeff., {\tt float} & Fr\'equence d'\'echantillonnage de $X$ ($Hz$)\\
$t_0$ & Coeff., {\tt float} & D\'ecalage ($s$)\\
\end{tabarg}

\medskip
%%%%%SORTIES%%%%
\begin{tabarg}
\multicolumn{3}{l}{\titretab {Sorties}}\\
\hline
$X_d$ & Vecteur, {\tt float} & Param\`etre d\'ecal\'e (m\^eme unit\'e que $X$ 
)\\
\end{tabarg}


%%%%%FORMULE%%%%
\itemstyle{Formule}

%%Decale
\begin{eqnarray*}
I_0 & = & Partie Entiere(-t_0 . f_e ) \\
F_0 & = & Partie Decimale(-t_0 . f_e ) \\
Xd[i] & = & X[i+I_0] + F_0.\left( X[i+I_0+1] + X[i+I_0] \right) \\
\end{eqnarray*}


\itemstyle{Commentaire~:}

Si $t_0$ est multiple de la p\'eriode
d'\'echantillonnage ($1/f_e$), alors le signal
est d\'ecal\'e d'un nombre entier de points, mais il
reste le m\^eme. Dans le cas contraire, 
le signal r\'esultat est interpol\'e 
lin\'eairement. Donc,  
la partie haute-fr\'equence est att\'enu\'ee. 

%%%%%AUTEUR%%%%
\itemstyle{Auteur~:}
CNRM/GMEI/TRAMM

%%%%%REFERENCE%%%%
\itemstyle{R\'ef\'erences~:}
